% =====================================================================
%  Laura P. Schaposnik
%  "Advanced topics in gauge theory: mathematics and physics of Higgs bundles"
%  arXiv:1907.09800v1 (PCMI 2019 Graduate Summer School 讲义, 53 页)
%  —— 中文译注本
%
%  编译：LuaLaTeX（本机统一用 lualatex，编译 3 遍）——见 build.ps1
%  约定：
%   * 原文公式编号在左侧（leqno），与原文版式一致；全部用 \tag{} 复刻原编号；
%   * 原文的定理/命题/例/注/猜想/问题与公式共用「节.序号」计数器（如 例 2.10、式 (2.12)），
%     本译本所有编号一律手工给出，绝不自动编号；
%   * % [Tutor Note: ...] 为译者批注，不进入排版，汇总见 译注汇总.md。
% =====================================================================
\documentclass[11pt,a4paper,leqno]{article}

\usepackage[UTF8]{ctex}
\usepackage{amsmath,amssymb,amsthm,mathtools}
\usepackage{bm}
\usepackage{graphicx}
\usepackage[a4paper,left=2.6cm,right=2.6cm,top=2.6cm,bottom=2.7cm]{geometry}
\usepackage[dvipsnames]{xcolor}
\usepackage{enumitem}
\usepackage{booktabs}
\usepackage{array}
\usepackage{multirow}
\usepackage{longtable}
\usepackage{multicol}
\usepackage{fancyhdr}
\usepackage{caption}
\usepackage{tikz}
\usepackage[hidelinks,unicode]{hyperref}
\hypersetup{pdftitle={规范场论中的进阶专题：Higgs 丛的数学与物理（中译注本）},
  pdfauthor={Laura P. Schaposnik; 中译/批注: 本译本}}

\captionsetup{font=small,labelsep=quad}
\setlength{\emergencystretch}{3em}
\allowdisplaybreaks[1]
\graphicspath{{figures/}{../figures/}}

% ---------- 页眉 ----------
\pagestyle{fancy}
\fancyhf{}
\fancyhead[L]{\small\itshape Advanced topics in gauge theory: mathematics and physics of Higgs bundles}
\fancyhead[R]{\small\thepage}
\renewcommand{\headrulewidth}{0.4pt}

% ---------- 编号条目环境（编号一律手工给定，见 STYLE_GUIDE §3） ----------
% 用法：\begin{thm}{2.6} …… \end{thm}
% 带原文括注时把括注写进参数：\begin{thm}{4.2（[59] 定理 11）}
\newcommand{\stmtopen}[1]{\par\medskip\noindent\textbf{#1}\ \kaishu}
\newenvironment{thm}[1]{\stmtopen{定理 #1.}}{\par\medskip}
\newenvironment{prop}[1]{\stmtopen{命题 #1.}}{\par\medskip}
\newenvironment{lem}[1]{\stmtopen{引理 #1.}}{\par\medskip}
\newenvironment{coro}[1]{\stmtopen{推论 #1.}}{\par\medskip}
\newenvironment{defn}[1]{\stmtopen{定义 #1.}}{\par\medskip}
\newenvironment{exmp}[1]{\stmtopen{例 #1.}}{\par\medskip}
\newenvironment{rem}[1]{\stmtopen{注 #1.}}{\par\medskip}
\newenvironment{conj}[1]{\stmtopen{猜想 #1.}}{\par\medskip}
% 习题集（原文 Problem set I–IV）中的「Problem N. 标题.」：正文为正常字体
\newenvironment{prob}[1]{\par\medskip\noindent\textbf{问题 #1.}\ }{\par\medskip}

% 原文作者用粗体行内小标题（如 "Bibliography."）——中文译为 \para{...}
\newcommand{\para}[1]{\par\medskip\noindent\textbf{#1}\quad\ignorespaces}

% 原文表格标题（手工编号）
\newcommand{\tabtitle}[2]{\caption*{\small\textbf{表 #1：}#2}}
% 原文图标题（手工编号）
\newcommand{\figtitle}[2]{\caption*{\small\textbf{图 #1：}#2}}

% 参考文献条目（gen_refs.py 按 \refitem{编号}{内容} 输出）
\newcommand{\refitem}[2]{\par\hangindent=1.5em\hangafter=1\noindent [#1]~#2}

% 习题集小节中罗马数字子项：\begin{enumerate}[label=(\roman*)]
% 原文 "i. ii. iii." 的缩进由 enumitem 控制（见 STYLE_GUIDE）

% ---------- 常用记号（只能用这里定义过的宏，不要自造 \newcommand） ----------
\newcommand{\R}{\mathbb{R}}
\newcommand{\C}{\mathbb{C}}
\newcommand{\Z}{\mathbb{Z}}
\newcommand{\N}{\mathbb{N}}
\newcommand{\Q}{\mathbb{Q}}
\newcommand{\E}{\mathbb{E}}
\newcommand{\PP}{\mathbb{P}}
\newcommand{\HH}{\mathbb{H}}
\newcommand{\T}{\mathbb{T}}
\newcommand{\cA}{\mathcal{A}}
\newcommand{\cB}{\mathcal{B}}
\newcommand{\cC}{\mathcal{C}}
\newcommand{\cD}{\mathcal{D}}
\newcommand{\cE}{\mathcal{E}}
\newcommand{\cF}{\mathcal{F}}
\newcommand{\cG}{\mathcal{G}}
\newcommand{\cH}{\mathcal{H}}
\newcommand{\cI}{\mathcal{I}}
\newcommand{\cJ}{\mathcal{J}}
\newcommand{\cK}{\mathcal{K}}
\newcommand{\cL}{\mathcal{L}}
\newcommand{\cM}{\mathcal{M}}
\newcommand{\cN}{\mathcal{N}}
\newcommand{\cO}{\mathcal{O}}
\newcommand{\cP}{\mathcal{P}}
\newcommand{\cQ}{\mathcal{Q}}
\newcommand{\cR}{\mathcal{R}}
\newcommand{\cS}{\mathcal{S}}
\newcommand{\cT}{\mathcal{T}}
\newcommand{\cU}{\mathcal{U}}
\newcommand{\cV}{\mathcal{V}}
\newcommand{\cW}{\mathcal{W}}
\newcommand{\cX}{\mathcal{X}}
\newcommand{\cY}{\mathcal{Y}}
\newcommand{\cZ}{\mathcal{Z}}
\newcommand{\fg}{\mathfrak{g}}
\newcommand{\fh}{\mathfrak{h}}
\newcommand{\fm}{\mathfrak{m}}
\newcommand{\fsl}{\mathfrak{sl}}
\newcommand{\fsp}{\mathfrak{sp}}
\newcommand{\fso}{\mathfrak{so}}
\newcommand{\fsu}{\mathfrak{su}}
\newcommand{\fu}{\mathfrak{u}}
\newcommand{\Gl}{\mathrm{GL}}
\newcommand{\Sl}{\mathrm{SL}}
\newcommand{\SO}{\mathrm{SO}}
\newcommand{\Spin}{\mathrm{Spin}}
\newcommand{\SU}{\mathrm{SU}}
\newcommand{\Sp}{\mathrm{Sp}}
\newcommand{\PSL}{\mathrm{PSL}}
\newcommand{\PGL}{\mathrm{PGL}}
\newcommand{\U}{\mathrm{U}}
\newcommand{\Prym}{\operatorname{Prym}}
\newcommand{\Rep}{\operatorname{Rep}}
\newcommand{\Hom}{\operatorname{Hom}}
\newcommand{\End}{\operatorname{End}}
\newcommand{\Sym}{\operatorname{Sym}}
\newcommand{\Jac}{\operatorname{Jac}}
\newcommand{\Pic}{\operatorname{Pic}}
\newcommand{\rk}{\operatorname{rk}}
\newcommand{\rank}{\operatorname{rank}}
\newcommand{\Tr}{\operatorname{Tr}}
\newcommand{\tr}{\operatorname{tr}}
\newcommand{\diag}{\operatorname{diag}}
\newcommand{\id}{\mathrm{id}}
\newcommand{\im}{\operatorname{im}}
\newcommand{\coker}{\operatorname{coker}}
\newcommand{\Spec}{\operatorname{Spec}}
\newcommand{\const}{\mathrm{const}}
\newcommand{\dd}{\mathrm{d}}
\newcommand{\ee}{\mathrm{e}}
\newcommand{\ii}{\mathrm{i}}
\newcommand{\ie}{\textit{i.e.}}
\newcommand{\eg}{\textit{e.g.}}
\newcommand{\cf}{\textit{cf.}}
\newcommand{\half}{\tfrac{1}{2}}

% ---------- 译者常用别名（2026-11-04 依译者实际用法补充） ----------
\newcommand{\SL}{\mathrm{SL}}
\newcommand{\GL}{\mathrm{GL}}
\newcommand{\GC}{G_{\C}}
\newcommand{\HC}{H_{\C}}
\newcommand{\Aut}{\operatorname{Aut}}
\newcommand{\Gr}{\operatorname{Gr}}
\newcommand{\Crit}{\operatorname{Crit}}
